\subsection{PRODUCT DECOMPOSITION OF THE AFFINE SPACE E}

We retain the notation of Prop. 5. For \(0 \leqslant p \leqslant s\) , let \(\mathbf{E}_p\) be the set of orbits of the group \(\mathbf{T}_p'\) in \(\mathbf{E}\) , and let \(\varphi_p\) be the canonical map from \(\mathbf{E}\) to \(\mathbf{E}_p\) . The action of \(\mathbf{T}\) on \(\mathbf{E}\) passes to the quotient; in particular, \(\mathbf{T}_p\) acts on \(\mathbf{E}_p\) and it is immediate (for example by taking an origin in \(\mathbf{E}\) ) that \(\mathbf{E}_p\) is an affine space admitting \(\mathbf{T}_p\) as its space of translations. Put \(\mathbf{E}' = \mathbf{E}_0 \times \cdots \times \mathbf{E}_s\) ; this is an affine space having \(\mathbf{T} = \mathbf{T}_0 \oplus \cdots \oplus \mathbf{T}_s\) as its space of translations. Let \(\varphi : \mathbf{E} \to \mathbf{E}'\) be the product map of the \(\varphi_p\) ; since \(\varphi\) commutes with the action of \(\mathbf{T}\) , this is a bijection and even an isomorphism of affine spaces. In what follows, we identify \(\mathbf{E}\) and \(\mathbf{E}'\) by means of \(\varphi\) ; the map \(\varphi_p\) is then identified with the canonical projection of \(\mathbf{E}'\) onto \(\mathbf{E}_p\) .

Since W leaves \(T_{p}^{\prime}\) stable, the action of W on E passes to the quotient and defines an action of W on \(E_{p}\) (for \(0 \leqslant p \leqslant s\) ), and hence by restriction an action of \(W_{p}\) on \(E_{p}\) (for \(1 \leqslant p \leqslant s\) ). On the other hand, let C be a chamber, let \(i \in I\) and let p be the integer such that \(i \in J_{p}\) . For any \(x \in E\) , we have

\[
s _ {i} (\mathrm{C}) (x) = x - \lambda . e _ {i} (\mathrm{C}) \quad \text { with } \quad \lambda \in \mathbf {R}.
\]

Since \(e_i \in \mathbf{T}_q'\) for \(q \neq p\) , it follows that

\[
\varphi_ {q} (w (x)) = \varphi_ {q} (x) \text {for} x \in \mathrm{E}, w \in \mathrm{W} _ {p}; 0 \leqslant q \leqslant s \text {and} q \neq p.
\]

Consequently, if \(w = w_{1} \ldots w_{s} \in \mathrm{W}\) with \(w_{p} \in \mathrm{W}_{p}\) for \(1 \leqslant p \leqslant s\) , then

\[
w \left(\left(x _ {0}, \dots , x _ {s}\right)\right) = \left(x _ {0}, w _ {1} \left(x _ {1}\right), \dots , w _ {s} \left(x _ {s}\right)\right)\tag{5}
\]

for all \(x_{p} \in E_{p}\) for \(0 \leqslant p \leqslant s\) . In other words, the action of W on \(E = E'\) is exactly the product of the actions of the \(W_{p}\) on \(E_{p}\) (we put \(W_{0} = \{1\}\) ). It follows that \(W_{p}\) acts faithfully on \(E_{p}\) and that \(W_{p}\) can be identified with a group of displacements of the euclidean space \(E_{p}\) (the space \(T_{p}\) of translations of \(E_{p}\) being provided, of course, with the scalar product induced by that on T).

PROPOSITION 6. (i) The group \(W_{p}\) is a group of displacements of the euclidean affine space \(E_{p}\) ; it is generated by reflections; provided with the discrete topology, it acts properly on \(E_{p}\) ; it is irreducible.

\begin{enumerate}
\def\labelenumi{(\roman{enumi})}
\setcounter{enumi}{1}
\tightlist
\item
  Let \(\mathfrak{H}_p\) be the set of hyperplanes \(\mathrm{H}\) of \(\mathrm{E}_p\) such that \(s_{\mathrm{H}} \in \mathrm{W}_p\) . The set \(\mathfrak{H}\) consists of the hyperplanes of the form
\end{enumerate}

\[
\mathrm{H} = \mathrm{E} _ {0} \times \mathrm{E} _ {1} \times \dots \times \mathrm{E} _ {p - 1} \times \mathrm{H} _ {p} \times \mathrm{E} _ {p + 1} \times \dots \times \mathrm{E} _ {s},
\]

with \(p = 1,\dots ,s\) and \(\mathrm{H}_p\in \mathfrak{H}_p\)

\begin{enumerate}
\def\labelenumi{(\roman{enumi})}
\setcounter{enumi}{2}
\tightlist
\item
  Every chamber C is of the form \(E_{0} \times C_{1} \times \cdots \times C_{s}\) , where for each p the set \(C_{p}\) is a chamber defined in \(E_{p}\) by the set of hyperplanes \(H_{p}\) ; moreover, the walls of \(C_{p}\) are the hyperplanes \(\varphi_{p}(H_{i}(C))\) for \(i \in J_{p}\) .
\end{enumerate}

Let \(\mathbf{C}\) be a chamber. Put \(\mathrm{H}_i = \mathrm{H}_i(\mathrm{C})\) , \(e_i = e_i(\mathrm{C})\) and \(s_i = s_i(\mathrm{C})\) for \(i \in \mathbb{I}\) (in the notation of no. 4).

\begin{enumerate}
\def\labelenumi{(\roman{enumi})}
\tightlist
\item
  Let \(i\) be in \(\mathbf{J}_p\) ; since \(e_i \in \mathrm{T}_p\) and \(\mathrm{T}\) is the direct sum of the mutually orthogonal subspaces \(\mathrm{T}_0, \mathrm{T}_1, \ldots, \mathrm{T}_s\) , the hyperplane of \(\mathrm{T}\) orthogonal to \(e_i\) is of the form \(\mathrm{L}_i + \mathrm{T}_p'\) , where \(\mathrm{L}_i\) is the hyperplane of \(\mathrm{T}_p\) orthogonal to \(e_i\) . The affine hyperplane \(\mathrm{H}_i\) of \(\mathrm{E}\) is of the form \(\mathrm{L}_i + \mathrm{T}_p' + x\) , with \(x \in \mathrm{E}\) , and we have
\end{enumerate}

\[
\mathrm{H} _ {i} = \mathrm{E} _ {0} \times \mathrm{E} _ {1} \times \dots \times \mathrm{E} _ {p - 1} \times \mathrm{H} _ {i} ^ {\prime} \times \mathrm{E} _ {p + 1} \times \dots \times \mathrm{E} _ {s}\tag{6}
\]

with \(H_{i}^{\prime}=L_{i}+\varphi_{p}(x)=\varphi_{p}(H_{i})\) . It is now immediate that \(s_{i}\) acts in \(E_{p}\) by the reflection associated to the hyperplane \(H_{i}^{\prime}\) of \(E_{p}\) . Thus, the group \(W_{p}\) is a group of displacements generated by reflections in \(E_{p}\) ; the verification of the properness criterion (D'2) is immediate. Finally, Prop. 5, (v) shows that \(W_{p}\) is irreducible. This proves (i).

\begin{enumerate}
\def\labelenumi{(\roman{enumi})}
\setcounter{enumi}{1}
\tightlist
\item
  By the Cor. of Th. 1, the set \(\mathfrak{H}_p\) consists of the hyperplanes of the form \(w_p(\mathrm{H}_i')\) for \(i\) in \(\mathrm{J}_p\) and \(w_p\) in \(\mathrm{W}_p\) . Further, if \(w = w_1 \ldots w_s\) with \(w_p \in \mathrm{W}_p\) for all \(p\) , formulas (5) and (6) imply that
\end{enumerate}

\[
w (\mathrm{H} _ {i}) = \mathrm{E} _ {0} \times \mathrm{E} _ {1} \times \dots \times \mathrm{E} _ {p - 1} \times w _ {p} (\mathrm{H} _ {i} ^ {\prime}) \times \mathrm{E} _ {p + 1} \times \dots \times \mathrm{E} _ {s},\tag{7}
\]

from which (ii) follows immediately.

\begin{enumerate}
\def\labelenumi{(\roman{enumi})}
\setcounter{enumi}{2}
\tightlist
\item
  Let \(i\) be in \(\mathbf{J}_p\) ; by formula (6), the open half-space \(D_i\) bounded by \(\mathbf{H}_i\) and containing \(C\) is of the form
\end{enumerate}

\[
\mathrm{D} _ {i} = \mathrm{E} _ {0} \times \mathrm{E} _ {1} \times \dots \times \mathrm{E} _ {p - 1} \times \mathrm{D} _ {i} ^ {\prime} \times \mathrm{E} _ {p + 1} \times \dots \times \mathrm{E} _ {s},
\]

where \(D_{i}^{\prime}\) is an open half-space bounded by \(H_{i}^{\prime}\) in \(E_{p}\) . Put \(C_{p} = \bigcap_{i \in J_{p}} D_{i}^{\prime}\) ; since \(C = \bigcap_{i \in I} D_{i}\) , it follows immediately that

\[
\mathrm{C} = \mathrm{E} _ {0} \times \mathrm{C} _ {1} \times \dots \times \mathrm{C} _ {s};
\]

consequently, none of the sets \(C_{p}\) is empty, and since C does not meet any hyperplane belonging to H, the set \(C_{p}\) does not meet any hyperplane belonging to \(H_{p}\) . Prop. 5 of §1, no. 3 now shows that \(C_{p}\) is one of the chambers defined by \(H_{p}\) in \(E_{p}\) . By using Prop. 4 of §1, no. 2, it is easy to see that the walls of \(C_{p}\) are the hyperplanes \(H_{i}^{\prime} = \varphi_{p}(H_{i})\) for \(i \in J_{p}\) .

